\chapter{Founder Psychology: Traits, Habits \& Expert Cognition}\label{ch:founder-psychology}

% Scope: the EVIDENCE-BASED psychology of the entrepreneur --- self-efficacy (Bandura +
%   Stajkovic & Luthans meta-analysis), goal-setting (Locke & Latham), personality traits
%   matched to the entrepreneurial task (Rauch & Frese), human capital outcomes vs.\ investments
%   (Unger et al.), deliberate practice with bounded-effect caveat (Ericsson; Charness;
%   Macnamara), effectuation (Sarasvathy), habit formation at 66 days (Lally), implementation
%   intentions (Gollwitzer & Sheeran), behavioral wealth-building (Thaler & Benartzi).
%   No debunk sections; only validated, actionable science.
% Cross-link heavily to \cref{ch:decisions-under-uncertainty} (biases) and
%   \cref{ch:organizational-behavior} (motivation/SDT) --- do NOT restate; \cref.

The entrepreneur-psychology literature is large, contentious, and often written backwards: it
starts with successful founders and reverse-engineers what they claim to have believed, done, or
felt. This chapter does the opposite. It starts from peer-reviewed meta-analyses and
pre-registered field studies, extracts the mechanisms and effect sizes that survive rigorous
scrutiny, and maps them onto the decisions a founder actually faces. The findings are more
modest than the bestseller covers imply, but they are more useful because they are
actionable: specific goals, domain-anchored confidence, traits matched to the entrepreneurial
task, expertise built through structured feedback, decisions calibrated to affordable loss, and
behavior wired into systems rather than left to willpower. These are the constructs with
replication behind them. They are also, as the sections below demonstrate, tightly wired to the
operating metrics of the running SaaS company that carries every worked example in this book.

\section{Self-Efficacy \& Goal-Setting}\label{sec:self-efficacy-goals}
% complexity: 2

What separates a founder who attempts a hard problem from one who merely discusses it? The
proximate answer is not talent, capital, or luck --- it is the belief that a specific course
of action, taken by them, will produce a result they can influence. Albert Bandura formalised
this construct as \emph{self-efficacy}: the judgment an individual makes about their capacity
to organise and execute the courses of action required by a specific situation
\parencite{bandura1997self}. Two features distinguish self-efficacy from vaguer concepts like
general confidence or optimism. First, it is domain-specific: a founder can have high
self-efficacy for product iteration and low self-efficacy for enterprise sales in the same
week. Second, it predicts behaviour via mechanism rather than correlation --- people high in
self-efficacy for a task approach obstacles as problems to be mastered, extend effort longer
when they encounter resistance, and recover faster after setbacks.

The meta-analytic evidence on this mechanism is unusually strong. Stajkovic and Luthans
synthesised \num{114} independent studies encompassing \num{21616} participants and found a
weighted average correlation between task-specific self-efficacy and work-related performance
of $r \approx .38$ % source: stajkovic1998self, Psychological Bulletin 1998, k=114, N=21,616
\parencite{stajkovic1998self}. In practical terms, moving from low to high self-efficacy in a
given domain is associated with roughly a \qty{28}{\percent} increase in output, operating
through three channels: high-efficacy individuals select more difficult and higher-yield goals,
exert greater initial effort when they begin a task, and persist longer when they encounter
resistance or market rejection. The channel matters as much as the correlation: self-efficacy
does not improve performance by making founders feel better; it improves performance by
changing the goals they set and the effort they sustain. Achievement motivation --- the intrinsic
drive to overcome challenges and meet high objective standards --- amplifies this further.
Collins, Hanges, and Locke's meta-analysis demonstrated that achievement motivation is
significantly correlated with both the choice of an entrepreneurial career and entrepreneurial
performance, with particularly strong differentiation when comparing high-performing founders
against lesser-performing peers % source: collins2004achievement, Human Performance 2004
\parencite{collins2004achievement}.

Self-efficacy answers \emph{whether} to act. Goal-setting answers \emph{where} to aim. Locke
and Latham's meta-analytic synthesis, built over thirty-five years of field and laboratory
studies, reached a deceptively simple conclusion: specific, difficult goals produce
systematically higher performance than vague exhortations to ``do your best,'' with effect
sizes ($d$) ranging from \num{0.52} to \num{0.82} across the synthesised trials
\parencite{locke2002building}. The mechanism runs through four channels: goals direct
attention to relevant activities and away from irrelevant ones; high goals raise the energy
applied to the task; difficult goals sustain effort longer because the gap between the current
state and the target keeps motivation active; and goals trigger strategy discovery --- when
existing approaches fall short, the goal creates pressure to search for better ones.

Two conditions gate the effect. First, goal commitment: if a target feels politically imposed
or genuinely impossible, the four-channel mechanism never engages. Second, task complexity:
for routine tasks, the ``specific and difficult'' rule is robust; for genuinely novel problems ---
a common situation in early-stage ventures --- rigid performance goals can produce panicked,
unsystematic behaviour. In those conditions, Locke and Latham recommend substituting a
\emph{learning goal} (``discover the three drivers of trial-to-paid conversion'') for a
performance goal (``reach \qty{25}{\percent} conversion this month''). The distinction matters
for founders because their environment mixes both: execution against known processes deserves
hard performance targets; genuine exploration deserves learning-goal framing.

\begin{example}[SaaS founder setting a milestone structure]
The SaaS founder running \money{340000} in monthly recurring revenue decides to lift the
trial-to-paid conversion rate, currently at \qty{25}{\percent}. She recognises that the channel
dynamics and attribution issues are already well-understood from the work in
\cref{ch:attribution-and-measurement}; the execution question is settled enough that a
performance goal is appropriate.

Following Locke and Latham, she sets a specific, difficult goal: reach \qty{30}{\percent}
trial-to-paid conversion within ninety days. She then decomposes it into weekly milestones ---
instrument the onboarding flow by day seven; run the first A/B test of the activation email by
day fourteen; assess cohort retention differences by day thirty. Each milestone has a
measurable output and a deadline, preserving the directive and energising functions of
goal-setting theory. Because the conversion process itself is well-understood (not novel), she
stays with performance framing throughout. She pairs the goal with a deliberate self-efficacy
check: having shipped three previous A/B experiments, she has domain-specific evidence that she
can execute this kind of project, which raises the probability that she will sustain effort when
the first test disappoints.

The implication: goal-setting works precisely because the goal structure --- specific, dated,
decomposed --- makes it impossible to substitute vague intention for concrete action. The
\qty{28}{\percent} performance lift associated with high task-specific self-efficacy compounds
with the goal-setting effect when both are present; neither operates well in the other's absence.
\end{example}

\begin{admonition}{Note}
Self-efficacy is built, not born. Bandura identifies four sources: mastery experiences (prior
successes in the same domain), vicarious modelling (watching comparable others succeed),
social persuasion (credible feedback from people who have observed the performance), and
physiological state (interpreting anxiety as arousal rather than incompetence). The practical
implication for founders is that self-efficacy in a new domain --- enterprise sales, financial
modelling, public speaking --- is raised most reliably by accumulating small, domain-specific
wins, not by general motivational input.
\end{admonition}

The intrinsic-motivation literature, surveyed in \cref{sec:motivation}, converges with
Bandura's model: autonomy and competence are the two self-determination-theory needs most
directly activated by a well-structured goal system. The goal gives the task a clear structure
(competence), while the founder's choice of the goal and method preserves ownership (autonomy).
\textcite{bandura1997self}, \textcite{stajkovic1998self}, and \textcite{locke2002building}
together supply the empirical ground; the practical question is always whether the task is
settled enough for a performance goal or novel enough to require a learning goal first.

\section{Traits \& Human Capital That Predict Venture Performance}\label{sec:traits-human-capital}
% complexity: 3

If self-efficacy and goal commitment explain how founders sustain effort, what personal
characteristics predict whether they will succeed at the task of business ownership in the first
place? For decades, entrepreneurship research produced inconsistent answers because it asked the
wrong question: whether broad personality traits --- extraversion, agreeableness, openness ---
predicted venture outcomes. They largely do not. The breakthrough came from reframing the
question: which traits are logically \emph{matched} to the specific task demands of running a
business, and which human capital dimensions are actually \emph{deployed} rather than merely
accumulated?

Rauch and Frese's meta-analysis provided the definitive quantification. Analysing personality
variables across a large set of independent samples, they found a stark contrast between traits
matched versus unmatched to the entrepreneurial task. Matched traits --- need for achievement,
stress tolerance, proactive personality, need for autonomy, innovativeness, and generalised
self-efficacy --- produced correlations with business creation of $r = .247$ ($K = 47$,
$N = \num{13280}$) and with business success of $r = .250$ ($K = 42$, $N = \num{5607}$).
Unmatched traits showed negligible predictive validity for success ($r = .028$, $K = 13$,
$N = \num{2777}$). % source: rauch2007person, European Journal of Work & Org Psychology 2007
The practical implication is that general extraversion or conscientiousness as abstract traits
are weak predictors; stress tolerance and proactive initiation --- traits structurally matched to
the demands of a volatile, high-stakes environment --- carry real signal
\parencite{rauch2007person}.

The human capital findings run parallel. Conventional business assumptions conflate two things
that are meaningfully distinct: the \emph{investments} an individual has made in human capital
(years of formal education, general management tenure, an MBA) versus the \emph{outcomes} those
investments may or may not have produced (specific, demonstrable, task-related knowledge and
skills). Unger, Rauch, Frese, and Rosenbusch synthesised \num{70} independent samples totalling
$N = \num{24733}$ businesses and found an overall relationship between general human capital and
venture success of $r_c = .098$ --- small in absolute terms
\parencite{unger2011human}. % source: unger2011human, Journal of Business Venturing 2011, N=24,733
But the aggregate masks the operationally critical finding: human capital \emph{outcomes}
predicted success at substantially higher rates than human capital \emph{investments}, and
task-related human capital vastly outperformed low task-related capital. A decade of general
management experience is a weak predictor; the specific ability to run a full-cycle enterprise
software sales process or architect a scalable data pipeline is a strong one. The reason is
direct: venture success is determined by the quality of decisions made in specific, high-stakes
contexts, and generic credentials signal little about performance in those contexts.

For founders, these two findings --- matched traits and outcome-oriented human capital --- have
concrete implications for self-assessment and early hiring. The trait that deserves attention is
not how extroverted or likeable a co-founder is, but whether they demonstrate the stress
tolerance and proactive orientation that Rauch and Frese identified as task-matched. The hiring
criterion that deserves attention is not credentials or years of experience, but the ability to
demonstrate task-relevant skill in a structured assessment.

\begin{example}[Trait and human-capital audit for a SaaS hiring decision]
The SaaS founder, at \money{340000} monthly recurring revenue with \money{3000000} in invested
capital and an \qty{11}{\percent} hurdle rate, is evaluating two candidates for a head-of-sales
role. Candidate A holds an MBA from a recognised institution and spent eight years in corporate
account management. Candidate B has a shorter general resume but has demonstrably closed six
enterprise SaaS contracts in the past two years at average contract values above \money{60000},
and handles the structured sales-call roleplay in the interview without avoidance.

The Unger et al.\ framework directs her to weight Candidate B's outcome-oriented human capital
--- specific, deployed, task-related skill --- over Candidate A's investment credentials. The
Rauch and Frese framework reinforces this: if Candidate B also scores higher on observable
stress tolerance and proactive task initiation during the interview process (responding to
rejection with a revised approach rather than withdrawal), those matched traits multiply the
human-capital advantage. The contribution margin at stake --- \money{21} per month per customer
converted by the sales hire --- gives the founder a quantitative frame for what the hiring
decision is worth. An additional hundred customers per year at \money{252} annual margin each,
discounted at the \qty{11}{\percent} hurdle rate as a growing perpetuity with \qty{3}{\percent}
growth, represents approximately \money{315000} of present value. A hiring decision calibrated
to outcome evidence rather than credential signalling is directly traceable to that number.

The implication: matched-trait and outcome-capital assessments are not soft add-ons to a hiring
process; they are the quantitatively best-supported predictors of whether the hire will produce
the metric outcomes the firm needs.
\end{example}

\begin{admonition}{Note}
The cognitive-bias literature is directly relevant to this domain. The overweighting of
credentials and general experience in hiring reflects both the representativeness heuristic
(``this person looks like successful people I know'') and availability bias (the successes of
well-credentialled leaders are more visible than the failures). The Rauch and Frese data are an
empirical antidote: the relevant question is not whether the candidate matches the prototype of
a successful leader, but whether their traits and demonstrated skills are matched to the
specific task demands of the role. The decision-theoretic scaffolding in
\cref{ch:decisions-under-uncertainty} develops this directly.
\end{admonition}

Both sets of findings connect to the value-creation spine: human capital outcomes are not a
fixed endowment but the cumulative product of deliberate, feedback-rich experience, which is
the subject of the following section. \textcite{rauch2007person} and \textcite{unger2011human}
supply the meta-analytic foundation; \textcite{collins2004achievement} connects the
achievement-motivation thread from \cref{sec:self-efficacy-goals} to the broader personality
picture.

\section{How Expertise Is Built \& How Founders Decide}\label{sec:practice-effectuation}
% complexity: 3

If matched traits and outcome-oriented human capital predict venture performance, the next
question is practical: how does a founder build those outcomes, and how should the resulting
expertise be deployed in the genuinely uncertain environment of an early-stage venture? The
answer comes from two bodies of research that are usually taught separately --- expertise
science and entrepreneurial cognition --- but are logically sequential.

Ericsson, Krampe, and Tesch-R{\"o}mer introduced deliberate practice to distinguish expert
performance from mere accumulated experience \parencite{ericsson1993role}. The key properties
are specificity, immediate feedback, and progressively increasing difficulty. Practice must
target a well-defined weakness; feedback on performance must be rapid and accurate; and the
challenge must be calibrated just above the current competence level. Work that falls within
established competence consolidates it but does not extend it. Ericsson spent the later part
of his career frustrated by the popular reduction of his findings to ``ten thousand hours of
anything,'' because the research shows that ten thousand hours of unstructured experience
produces a very good amateur, not an expert.

Field evidence from chess expertise confirms the mechanism. Charness, Tuffiash, Krampe,
Reingold, and Vasyukova demonstrated that solitary, serious study --- deliberate practice in the
technical sense --- was the single strongest predictor of skill level among players, accounting
for large variance in Elo ratings that accumulated experience alone could not explain
\parencite{charness2005practice}. % source: charness2005practice, Applied Cognitive Psychology 2005
The chess domain is instructive precisely because it has stable rules, clear objective
metrics, and rapid feedback loops --- the conditions under which deliberate practice is most
powerful.

Practitioners must apply a critical caveat when translating this to the business environment.
Macnamara, Hambrick, and Oswald's meta-analysis quantified the proportion of performance
variance explained by deliberate practice across multiple domains. % source: macnamara2014practice, Psychological Science 2014
In games with stable rules (including chess), deliberate practice explained \qty{26}{\percent}
of variance in performance. In music, it explained \qty{21}{\percent}. In less structured
professional domains with noisy, slow feedback, the figure falls substantially
\parencite{macnamara2014practice}. For founders, the implication is calibrated rather than
discouraging: deliberate practice is highly effective for isolated, well-defined sub-skills ---
objection handling in enterprise sales, financial modelling in a spreadsheet, writing a customer
discovery interview script --- but it is not the primary mechanism for ``company building'' as
a holistic, ambiguous activity. The Unger et al.\ human-capital finding (\cref{sec:traits-human-capital})
reinforces this: the human capital that predicts venture success is task-related and deployable,
not generic.

The entrepreneurial environment makes genuine deliberate practice difficult to construct at the
firm level because feedback loops are slow, noisy, and expensive. A surgeon knows within minutes
whether a procedure succeeded; an entrepreneur may wait twelve months to learn whether a product
decision was correct. The practical response is to create artificial feedback structures: rigorous
A/B testing provides the rapid, objective signal that a coach provides in a traditional expertise
domain; financial modelling exposes assumptions to falsification before the market does; advisory
relationships with experienced operators function as the expert coach the environment does not
automatically provide.

Expert entrepreneurial decision-making has a second distinguishing feature, identified by
Sarasvathy through a landmark study of thirty expert founders --- entrepreneurs with at least
fifteen years of experience who had taken at least one company public \parencite{sarasvathy2001causation}.
Given the same business case, expert founders did not reason the way MBA textbooks prescribe.
They did not start from a target market, project returns, and marshal resources toward the goal.
Instead, they started from their available means --- who they are, what they know, who they know
--- and allowed goals to emerge from what those means could plausibly produce. Losses were
bounded by an ``affordable loss'' threshold: not the maximum expected return, but the maximum
they were willing to put at risk. Sarasvathy called this pattern \emph{effectuation}, contrasted
with the \emph{causation} logic of standard strategic planning.

The distinction matters because it is not simply a stylistic preference. Expert founders use
effectuation because it is better adapted to environments of genuine uncertainty, where the
probability distributions required by expected-value reasoning do not exist. In that regime,
controlling the downside and exploiting contingencies as they arise dominates trying to predict
and optimise. The causal approach --- market research, projected return, resource acquisition ---
is appropriate when the future is sufficiently predictable and the market is sufficiently stable.
Early-stage entrepreneurship rarely satisfies either condition. This connects directly to the
decision-theoretic scaffolding in \cref{ch:decisions-under-uncertainty}: effectuation is, in
essence, the practitioner's implementation of robust decision-making under deep uncertainty.

\begin{example}[Effectual vs.\ causal framing for the SaaS]
The SaaS founder, currently at \money{340000} monthly recurring revenue with a
\qty{92.5}{\percent} twelve-month retention rate, is considering entering a new vertical ---
small law firms. A causal decision-maker would begin by estimating the total addressable market,
project a penetration rate, model the contribution margin from the new vertical (\money{21} per
month per customer, as established in the firm's cost structure), compute the net present value
of the expansion using the \qty{11}{\percent} hurdle rate, and acquire resources to execute if
NPV is positive.

An effectual founder asks different first questions: Do we know any partners at small law firms
who would try this for free and give us honest feedback? Can we configure the existing product
for the vertical within the current engineering sprint, or does it require a separate codebase?
What is the maximum time and budget we are willing to lose if this turns out to be a dead end ---
and is that amount genuinely affordable against current runway? She does not refuse to model
returns, but she treats the model as one input among several and weights the accessible evidence
from means --- relationships, existing capability, affordable exposure --- more heavily than the
projected evidence from a market forecast.

The implication: the two frameworks are not competitors for all situations. The causal approach
becomes appropriate once market response is better understood; effectuation is the right posture
when the uncertainty is genuine rather than resolvable through more research.
\end{example}

\begin{admonition}{Note}
Deliberate practice and effectuation are not in tension; they are sequential. Effectuation is
the decision logic for entering an uncertain situation with bounded risk. Deliberate practice is
what builds the domain competence that makes future effectuation richer --- because the founder's
means (what they know, what they can do) expand with each feedback-rich iteration. The expert
entrepreneurs in Sarasvathy's study had extensive deliberate-practice histories in their domains;
effectuation was how they leveraged that expertise in new situations, not a substitute for having
built it.
\end{admonition}

\section{Habits \& Behavioral Wealth-Building}\label{sec:habits-wealth}
% complexity: 2

The framing of high performance as a product of willpower and motivation is one of the most
reliably wrong assumptions in popular psychology. Willpower depletes; motivation fluctuates;
neither is a dependable foundation for the long-run behaviour change that compound growth ---
personal or financial --- requires. The behavioural science of habit formation and choice
architecture offers a more durable alternative: design the environment so that the desired
behaviour is the path of least resistance, and the undesired behaviour requires an active
decision to access.

The foundational scientific model for understanding how habits interact with goals is provided by
Wood and Neal, who demonstrated that habits emerge through the gradual learning of associations
between responses and specific features of performance contexts --- time of day, physical location,
or preceding action sequences \parencite{wood2007habits}. Once a habit is formed, the perception
of the contextual cue triggers the associated behavioural response without the mediation of an
active goal state. Goals are required to motivate the initial repetitions that establish the
habit, but once automaticity is achieved, the behaviour becomes substantially resistant to
changes in the individual's motivational state. For founders, this automaticity preserves
executive function for complex problem-solving --- the cognitive resource most at risk during the
high-stress phases of a venture.

The practical question that popular habit books do not answer rigorously is how long this
formation process takes. Charles Duhigg, James Clear, and B.\ J.\ Fogg have all described habit
architecture compellingly, and their frameworks are useful design guides at the structural level
\parencite{duhigg2012power,clear2018atomic,fogg2020tiny}. But the ``21-day habit'' claim
embedded in popular culture is not empirically grounded. Lally, van Jaarsveld, Potts, and Wardle
tracked ninety-six volunteers over eighty-four days as they attempted to adopt a new health
behaviour linked to a daily cue, mapping automaticity using the Self-Report Habit Index
\parencite{lally2010habit}. % source: lally2010habit, European Journal of Social Psychology 2010, N=96, 84 days
The automaticity curve was asymptotic: large gains in early repetitions, diminishing returns
thereafter, plateauing at an individual maximum. The average time required to reach
\qty{95}{\percent} of the automaticity asymptote was \textbf{66 days}, with individual variance
ranging from 18 to 254 days depending on the complexity of the behaviour. Crucially, missing a
single opportunity to perform the behaviour did not materially affect the formation process or
alter the final asymptote --- meaning rigid unbroken streaks are not neurologically necessary,
but a sustained sixty-plus-day horizon is.

Because habits take a median of two months to form, founders require a reliable bridge between
intention and automaticity. That bridge is the implementation intention. Gollwitzer and Sheeran
analysed this mechanism across ninety-four independent tests and found a robust, medium-to-large
effect on goal attainment ($d \approx .65$) \parencite{gollwitzer2006implementation}.
% source: gollwitzer2006implementation, Advances in Experimental Social Psychology 2006, 94 tests, d≈.65
An implementation intention is an explicit ``if-then'' plan that specifies when and where a
goal-directed behaviour will be initiated: ``If it is 8:00 AM on a weekday, then I will open the
MRR dashboard before any other application.'' This format transfers the trigger for action from
internal willpower to an external environmental cue, promoting rapid initiation of goal-directed
behaviour, shielding ongoing goal pursuit from competing demands, and reducing the cognitive load
required to begin. The key mechanism is specificity: the ``if-then'' format pre-commits the
response to a context cue, so that when the cue appears, the behaviour executes without
requiring a decision.

The behavioural-finance application extends the same logic to personal financial decisions.
Entrepreneurs face an acute version of the standard self-control problem: irregular income,
strong present-biased preferences, and the constant temptation to redeploy personal savings into
the business. Thaler and Benartzi's ``Save More Tomorrow'' programme provides the
evidence-based counter-strategy \parencite{thaler2004save}. The programme's key design insight
is that loss aversion --- the asymmetry by which losses loom larger than equivalent gains in
subjective experience --- makes immediate cuts to take-home pay painful even when the long-term
benefit is obvious. Save More Tomorrow sidesteps this by asking participants to pre-commit to
directing a fixed fraction of their \emph{future} raises to savings. Because the deduction
coincides with an increase in gross pay, take-home income never falls; there is no loss-aversion
trigger. Thaler and Benartzi tracked participants over forty months and observed average saving
rates rise from \qty{3.5}{\percent} to \qty{13.6}{\percent} --- a nearly four-fold improvement
achieved without willpower, through architecture alone. The mechanism is identical to the habit
and implementation-intention literature: remove the decision from the domain of in-the-moment
deliberation; encode it in advance as an automatic commitment.

\begin{example}[Implementation intentions and automated saving for the SaaS founder]
The SaaS founder wants to establish two reliable routines: a weekly financial review and a
personal wealth-building protocol. She applies the Lally et al.\ finding that habit
automaticity requires a sixty-six-day median horizon, and the Gollwitzer and Sheeran finding
that an explicit ``if-then'' format ($d \approx .65$) bridges the gap while the habit forms.

For the financial review, she writes the implementation intention explicitly: ``If it is Friday
morning and my calendar shows 9:00 AM, then I will open the MRR dashboard, check burn rate
against the \money{3000000} invested-capital base, and record one decision or updated hypothesis
before closing it.'' The cue (Friday 9:00 AM) is stable and external. The routine (three-item
sequence) is specific enough to execute without further deliberation. After ten weeks --- within
the Lally et al.\ range --- the review has become context-cued rather than motivation-dependent:
it happens whether or not she feels engaged with it.

On the wealth-building side, she applies Thaler and Benartzi's architecture directly: each
time the business makes a quarterly distribution, \qty{20}{\percent} routes automatically to a
personal investment account before it reaches her operating account. The instruction is locked
in the bank's standing-order system and requires deliberate action to cancel. The trigger is an
incoming-cash event, not a spending decision, so loss aversion does not engage. The equivalent
of pre-committing to future salary increases is pre-committing to future distribution fractions:
the sacrifice is prospective, not immediate.

The implication is general: both the professional habit and the savings rule work because they
are architecture, not resolve. The implementation-intention format handles the forty-to-sixty-
day bridge before automaticity; the choice-architecture lock handles the long-run default. The
moment either requires the founder to decide in real time, under competing pressures, they
become unreliable.
\end{example}

\begin{admonition}{Note}
The autonomy need identified in self-determination theory (\cref{sec:motivation}) is not
undermined by habit automation --- it is expressed at the moment of designing the system. The
founder who pre-commits to a saving schedule and writes an implementation intention for a weekly
financial review has exercised high autonomy in constructing the choice architecture; what is
automated afterward is the execution, not the choice. The distinction matters because founders
sometimes resist systematic approaches as constraints on their agency, when in fact designing
effective systems is one of the highest-leverage exercises of agency available to them.
\end{admonition}

The chapter's argument closes where it opened: the entrepreneurial psychology that survives
contact with replication is not a personality profile or a belief system. It is a set of
practices --- structured goals paired with domain-specific confidence built through evidence,
traits and human capital matched to the task demands, decision logic adapted to genuine
uncertainty, and behaviour encoded in systems rather than resolved in real time. These are the
tools. \textcite{bandura1997self}, \textcite{stajkovic1998self}, \textcite{locke2002building},
\textcite{rauch2007person}, \textcite{unger2011human}, \textcite{ericsson1993role},
\textcite{sarasvathy2001causation}, \textcite{lally2010habit}, \textcite{gollwitzer2006implementation},
and \textcite{thaler2004save} provide the empirical foundation; the frameworks in
\cref{ch:decisions-under-uncertainty} and \cref{sec:motivation} locate them within the broader
decision-making and motivation science the book has already developed.
